% ---------------------------------------------------------------
\section{Anchored Energy Matching}
\label{sec:method}
% ---------------------------------------------------------------

\paragraph{Setting.}
Let $\pi(x)=\tilde\pi(x)/Z$ with $\tilde\pi(x)=\exp(-\mathcal{E}(x))$ on $\mathbb{R}^d$.
We use the variance-exploding (VE) noising process $x_t = x_0 + \sigma_t\varepsilon$, $\varepsilon\sim\mathcal{N}(0,I)$, with marginals $p_t=\pi*\mathcal{N}(0,\sigma_t^2 I)$, and write
$\pi_t(x_0\mid x_t)\propto \tilde\pi(x_0)\,\mathcal{N}(x_t;x_0,\sigma_t^2I)$ for the denoising posterior.
For any test function $g$, define the posterior mean $\bar g_t(x_t)=\mathbb{E}_{\pi_t(\cdot\mid x_t)}[g(x_0)]$.
Two choices of $g$ recover the noised score,
\begin{align}
\nabla\log p_t(x_t)
&= \mathbb{E}_{\pi_t(\cdot\mid x_t)}\big[-\nabla\mathcal{E}(x_0)\big] && \text{(TSI)}\label{eq:tsi}\\
&= \mathbb{E}_{\pi_t(\cdot\mid x_t)}\big[(x_0-x_t)/\sigma_t^2\big] && \text{(DSI)}.\label{eq:dsi}
\end{align}
The target score identity (TSI) needs energy gradients. The denoising score identity (DSI) needs only energies, but it is useless without posterior samples.

\paragraph{Energy matching and its bias.}
iDEM~\citep{akhound2024iterated} trains $s_\theta(x,t)$ by regressing onto a self-normalized importance sampling (SNIS) estimate of \eqref{eq:tsi}. It uses $K$ proposals $y_{1:K}\overset{\text{iid}}{\sim} q_t(\cdot\mid x_t)=\mathcal{N}(x_t,\sigma_t^2I)$ and importance weights $w(y)=\tilde\pi(y)\mathcal{N}(x_t;y,\sigma_t^2I)/q_t(y\mid x_t)=\tilde\pi(y)$:
\begin{equation}
\widehat S_K(x_t)=\frac{\sum_{k=1}^K w(y_k)\,g(y_k)}{\sum_{k=1}^K w(y_k)}.
\label{eq:snis}
\end{equation}
The noised points come from a replay buffer $\mathcal{B}$ filled with the model's own samples: $x_0\sim\rho$ (the buffer law) and $x_t=x_0+\sigma_t\varepsilon$.
The loss is $\mathcal{L}_\rho(\theta)=\mathbb{E}_{t,x_0,\varepsilon,y}\,\lambda_t\|s_\theta(x_t,t)-\widehat S_K(x_t)\|^2$.
Its minimizer over all functions is the conditional mean $s^\star_\rho(x,t)=\mathbb{E}[\widehat S_K(x_t)\mid x_t=x]$.
Because $\widehat S_K$ never looks at $x_0$, this minimizer is the \emph{same biased field} $b_K(x,t)=\mathbb{E}_{q_t}[\widehat S_K(x)]$ for \emph{every} buffer law $\rho$, including the target $\rho=\pi$.
SNIS has bias of order $\mathrm{Var}_{q_t}(w)/(K\,\mathbb{E}_{q_t}[w]^2)$. This term grows exponentially as $\sigma_t$ grows, because the proposal spreads over the whole space while $\tilde\pi$ concentrates on its modes.
So the true score is not a stationary point of iDEM training, and the error sits at high noise, which is where the sampler decides how much mass each mode gets.

\paragraph{The anchor.}
Each training pair already contains a candidate posterior sample: the clean point $x_0$ that generated $x_t$.
If $x_0\sim\pi$, then $x_0\mid x_t\sim\pi_t(\cdot\mid x_t)$ \emph{exactly}.
Anchored Energy Matching (AEM) adds $x_0$ to the particle system as the anchor $y_0:=x_0$:
\begin{equation}
\widehat A_K(x_t,x_0)=\frac{\sum_{k=0}^K w(y_k)\,g(y_k)}{\sum_{k=0}^K w(y_k)},\qquad y_0=x_0.
\label{eq:aem}
\end{equation}
This costs one extra energy (and gradient) evaluation per training pair, often a cached one.
The construction is the Rao--Blackwellized form of one step of the iterated sampling--importance resampling (i-SIR) kernel, i.e.\ conditional SMC with one time step~\citep{andrieu2010particle,andrieu2018uniform}.
Here we use it to build an unbiased \emph{regression target}, not to run a Markov chain.

% ---------------------------------------------------------------
\subsection{Theory}
\label{sec:theory}
% ---------------------------------------------------------------
Fix $t$ and $x_t$. Write $\pi_t=\pi_t(\cdot\mid x_t)$, $q=q_t(\cdot\mid x_t)$, and let $\bar w(y)=\pi_t(y)/q(y)$ be the normalized weight.
Assume $M_t(x_t):=\sup_y \bar w(y)<\infty$. This holds for the Gaussian proposal whenever $\tilde\pi$ is bounded.

\begin{theorem}[Exact unbiasedness and Rao--Blackwell variance bound]
\label{thm:unbiased}
Let $y_0\sim\pi_t$ and $y_{1:K}\overset{\text{iid}}{\sim}q$ be independent. For every $K\ge 0$ and every $g\in L^2(\pi_t)$,
\[
\mathbb{E}\big[\widehat A_K\big]=\bar g_t(x_t)
\quad\text{and}\quad
\mathrm{Cov}\big[\widehat A_K\big]\preceq \mathrm{Cov}_{\pi_t}\big[g\big].
\]
\end{theorem}
\begin{proof}[Proof sketch]
Pick an index $J$ uniformly from $\{0,\dots,K\}$ and swap the labels of $y_0$ and $y_J$.
This does not change the law of $\widehat A_K$, because the estimator is symmetric in its particles.
The relabelled system has joint law $\frac{1}{K+1}\pi_t(y_J)\prod_{k\neq J}q(y_k)$.
Summing over $J$ gives the marginal $\prod_k q(y_k)\cdot\frac{1}{K+1}\sum_j\bar w(y_j)$.
So the conditional law of $J$ given the particles is $\bar w(y_J)/\sum_j \bar w(y_j)$, while $y_J\sim\pi_t$ marginally.
Hence $\widehat A_K=\mathbb{E}[g(y_J)\mid y_{0:K}]$, and the tower rule gives $\mathbb{E}\widehat A_K=\mathbb{E}_{\pi_t}g$.
The covariance bound is the law of total covariance (Rao--Blackwell). Full proof in \cref{app:proofs}.
\end{proof}

Two consequences follow. First, $K=0$ recovers target score matching ($g=-\nabla\mathcal{E}$) or plain denoising score matching ($g$ from DSI); as $K\to\infty$ the anchor weight goes to zero and AEM tends to iDEM. AEM therefore interpolates between data-driven and energy-driven score learning, and every member of the family is unbiased.
Second, the DSI choice of $g$ becomes viable. Its variance $\mathrm{tr}\,\mathrm{Cov}_{\pi_t}[x_0]/\sigma_t^4$ decays as $\sigma_t$ grows, whereas the TSI variance $\mathrm{tr}\,\mathrm{Cov}_{\pi_t}[\nabla\mathcal{E}]$ does not, which explains the large gains at high noise in \cref{sec:exp-estimator}.

\begin{corollary}[The target is a fixed point]
\label{cor:fixed}
Let $s^{\mathrm{A}}_\rho(x,t)=\mathbb{E}[\widehat A_K\mid x_t=x]$ be the population minimizer of the AEM loss with buffer law $\rho$.
Then $s^{\mathrm{A}}_\pi=\nabla\log p_t$ for all $K\ge0$.
In contrast, the iDEM minimizer is $b_K$ for every $\rho$, and $b_K\neq\nabla\log p_t$ in general (\cref{app:idem-bias} gives a closed-form Gaussian counterexample).
\end{corollary}

In practice the buffer is only approximately $\pi$. The next result shows that AEM \emph{shrinks} the buffer error instead of copying it.

\begin{theorem}[Bias contraction under an approximate anchor]
\label{thm:contraction}
Let $y_0\sim\mu$ for an arbitrary distribution $\mu$ (e.g.\ $\mu=\rho_t(\cdot\mid x_t)$, the buffer posterior), with the other particles as before.
Then $\mathbb{E}[\widehat A_K]=(\mu P_K)(g)$, where $P_K$ is the i-SIR Markov kernel, which leaves $\pi_t$ invariant. For bounded $g$,
\[
\big\|\mathbb{E}[\widehat A_K]-\bar g_t(x_t)\big\|
\le 2\|g\|_\infty\,(1-\epsilon_K)\,\mathrm{TV}(\mu,\pi_t),
\qquad
\epsilon_K=\frac{K}{K-1+2M_t}.
\]
Running $n$ i-SIR steps before the final Rao--Blackwellized average (each step resamples a new anchor and redraws the proposals) replaces $(1-\epsilon_K)$ by $(1-\epsilon_K)^n$.
\end{theorem}
\begin{proof}[Proof sketch]
The first claim is the identity $\widehat A_K=\mathbb{E}[g(y_J)\mid y_{0:K}]$ with $J\sim\mathrm{Cat}(\bar w)$, which holds for any anchor law.
For the minorization, any $y_0$ and measurable $B$ satisfy
$P_K(y_0,B)\ge K\,\mathbb{E}\big[\bar w(y_1)\mathbf 1_B(y_1)/(\bar w(y_0)+\bar w(y_1)+\sum_{k\ge2}\bar w(y_k))\big]\ge \epsilon_K\,\pi_t(B)$,
using $\bar w\le M_t$, the independence of $y_1$ from $y_{2:K}$, Jensen's inequality for $u\mapsto 1/(2M_t+u)$, and $\mathbb{E}_q\bar w=1$.
A Doeblin minorization with constant $\epsilon_K$ contracts total variation by $(1-\epsilon_K)$ per step.
\end{proof}

\paragraph{Reading the bound.}
\Cref{thm:contraction} is a pointwise statement: for each $x_t$, the regression target is closer to the true posterior mean than the buffer's posterior is to the true posterior, by a factor $1-\epsilon_K$. It does not by itself prove that the outer loop converges. That would also require controlling how the learned score maps back into a new buffer, which we leave open.
It also explains graceful degradation. A poor anchor (low $\tilde\pi(x_0)$) gets negligible weight and AEM reduces to iDEM. A good anchor dominates exactly when the proposals miss every mode, which is where iDEM's SNIS collapses onto whichever proposal lands closest to a mode.
The constant is also honest about the limitation. For Gaussian proposals, $M_t$ grows exponentially with dimension (for a unit Gaussian component near its mode, $M_t\approx e^{d/2}$). So for $d\ge16$ and practical $K$, the guaranteed contraction is weak ($\epsilon_K\ll1$), and $1-\epsilon_K$ decays only like $1/K$. In that regime the anchored target leans on the buffer, and the theorem predicts that AEM keeps, rather than corrects, buffer errors that no proposal can see. \Cref{sec:exp-ablation} shows exactly this.
Finally, the TSI and DSI test functions are unbounded. For them the bound holds after clipping $g$, which we do in practice, or in a weighted norm under a drift condition that we do not verify.

\subsection{Replay proposals: letting the energy correct the buffer}
\label{sec:replay}
\Cref{thm:unbiased,thm:contraction} hold for \emph{any} proposal $q$ that is independent of the anchor and has a density we can evaluate. The Gaussian proposal $\mathcal{N}(x_t,\sigma_t^2I)$ ignores the prior factor $\tilde\pi$ of the posterior. At high noise this makes $M_t$ huge, so contraction is weak and the anchor dominates.
A replay buffer gives us something better. For row $i$ of a minibatch with anchors $x_0^{(1:B)}$, define the leave-one-out mixture
\begin{equation}
q_i(y)=\alpha\,\mathcal{N}(y;x_t^{(i)},\sigma_t^2I)+(1-\alpha)\!\sum_{j\neq i}r_{ij}\,\mathcal{N}(y;x_0^{(j)},\tau_t^2I),
\label{eq:replay}
\end{equation}
with responsibilities $r_{ij}\propto\mathcal{N}(x_t^{(i)};x_0^{(j)},(\sigma_t^2+\tau_t^2)I)$ and bandwidth $\tau_t=\min(\sigma_t,\tau_{\max})$.
Because $q_i$ excludes row $i$, it is independent of the anchor $x_0^{(i)}$, as long as minibatch rows are distinct buffer points, and its density is available in closed form.\footnote{The training runs in \cref{sec:experiments} drew minibatch rows with replacement. The same buffer point can then appear in another row (probability $\approx2.5\%$ per row for $B=256$ and a buffer of $10^4$), which slightly violates independence. Our code provides sampling without replacement, which we use for the oracle experiment in \cref{sec:exp-ablation}.}
The weights $w(y)=\tilde\pi(y)\mathcal{N}(x_t;y,\sigma_t^2I)/q_i(y)$ are therefore exact, and both theorems apply unchanged.
The mixture has three effects.
(i)~Proposals now land in the modes the buffer has already found, so $M_t$, and with it the bias constant $1-\epsilon_K$, shrinks sharply.
(ii)~The importance weights compare energies \emph{across} modes. If the buffer over-represents a mode, $q_i$ is large there and $w$ is small, so the target corrects the mode weights; an anchor alone cannot do this.
(iii)~The Gaussian component ($\alpha>0$) keeps iDEM's exploration of modes the buffer has not found.
The cost is one $B\times B$ distance matrix and a $B\times K\times B$ log-sum-exp per step, which is negligible next to the energy evaluations.

\paragraph{Test function and practical deviations.}
We combine the TSI and DSI test functions with the fixed weights $\alpha_t=(1+\sigma_t^2)^{-1}$ for every method; any such combination keeps the target unbiased (\cref{app:extra}). The implementation also clips gradients and caps $\|\sigma_t\,\mathrm{target}\|$ identically for all methods, which introduces a small bias, and the finite buffer is only approximately $\pi$ (\cref{app:extra}).

\subsection{Refining the buffer: an exact i-SIR Gibbs step}
\label{sec:refine}
Anchoring leaves the buffer untouched, so it cannot repair buffer errors that the proposals do not see. We add an optional \emph{refinement} step that moves the buffer towards $\pi$ using only energy evaluations, and that leaves an exact buffer exact.
Treat the buffer as $N$ chains $x^{1:N}$ and a block $S$ of them as the state to be updated, with the remaining chains $R$ held fixed as reference.
For $x=x^i$, $i\in S$, draw a noise level $\sigma$ and apply two steps:
(1)~\emph{noise}: $x_t=x+\sigma\varepsilon$, an exact draw from $\mathcal{N}(x_t\mid x,\sigma^2I)$;
(2)~\emph{i-SIR resample}: draw $K$ proposals $y_k\sim q(\cdot\mid x_t,R)$, where $q$ is the mixture \eqref{eq:replay} built from the reference chains, and select $x'$ among $\{x,y_1,\dots,y_K\}$ with probability proportional to $\tilde\pi(y)\,\mathcal{N}(x_t;y,\sigma^2I)/q(y\mid x_t,R)$.
Step~(2) is the conditional importance sampling kernel of \cref{sec:method} applied to the posterior $\pi_t(\cdot\mid x_t)$, with the current chain playing the role of the anchor. We call the pair a refinement sweep.

\begin{proposition}[Invariance of refinement]
\label{prop:refine}
Let $\mathcal{P}_R$ be the kernel that updates every chain in $S$ independently by steps (1)--(2), given reference chains $R$. If $x^{S}\sim\pi^{\otimes|S|}$ then $\mathcal{P}_R$ maps $x^S$ to a sample of $\pi^{\otimes|S|}$, for every fixed $R$ and every $K\ge1$. Consequently, a sweep that visits all blocks in turn, each with a reference set made of chains from outside the block, leaves $\pi^{\otimes N}$ invariant.
\end{proposition}

\noindent The proof is in \cref{app:proofs} (it is the standard Gibbs-plus-i-SIR argument on the joint density $\pi(x)\mathcal{N}(x_t\mid x,\sigma^2I)$). Two features matter. The proposal depends on $R$ and $x_t$ but not on $x$, which is what makes step~(2) a valid i-SIR kernel. And the invariance is for the \emph{exact} buffer: a buffer that is not exact is not repaired in one step, and we prove no rate for it. Unlike anchoring, however, the step can change the relative weights of the modes the buffer contains, because the weights $\tilde\pi/q$ compare energies across modes.
Refinement requires $K{+}1$ energy evaluations per buffer point and sweep, with no gradients. It is run once after each outer iteration, once anchoring has started (explore$\to$anchor), with noise levels log-uniform in $[\sigma_{\min},\sigma_{\mathrm{ref}}]$. Since $\sigma_{\mathrm{ref}}$ sets how far a refinement step can move a point, we take it at the within-mode scale of the target (\cref{app:exp}). The proposals are the replay mixture with $\alpha=0.5$.

\begin{algorithm}[t]
\caption{Anchored Energy Matching (AEM)}
\label{alg:aem}
\begin{algorithmic}
\STATE {\bfseries Input:} energy $\mathcal{E}$, network $s_\theta$, buffer $\mathcal{B}$, proposals $K$, i-SIR steps $n$
\REPEAT
\STATE \emph{Outer:} simulate the reverse SDE with $s_\theta$; push samples into $\mathcal{B}$
\FOR{inner step $=1$ {\bfseries to} $N_{\mathrm{in}}$}
\STATE $x_0\sim\mathcal{B}$, $t\sim\mathcal{U}[0,1]$, $x_t=x_0+\sigma_t\varepsilon$
\STATE $y_0\gets x_0$ \hfill\COMMENT{anchor}
\FOR{$i=1$ {\bfseries to} $n$}
\STATE $y_{1:K}\sim\mathcal{N}(x_t,\sigma_t^2I)$; $\;\omega_k\propto \exp(-\mathcal{E}(y_k))$, $k=0..K$
\STATE {\bfseries if} $i<n$: $y_0\gets y_J$, $J\sim\mathrm{Cat}(\omega)$
\ENDFOR
\STATE $\widehat A\gets\sum_{k=0}^K\omega_k\,g_t(y_k)$
\STATE $\theta\gets\theta-\eta\nabla_\theta\,\lambda_t\|s_\theta(x_t,t)-\widehat A\|^2$
\ENDFOR
\UNTIL{budget exhausted}
\end{algorithmic}
\end{algorithm}
